% Copyright (c) 2014,2016,2018 Casper Ti. Vector
% Public domain.

\chapter{引言}

近年来，\ac{LLM} 与 \ac{VLM} 通过规模扩增与训练范式革新，在自然语言理解、逻辑推理与多模态感知方面展现出卓越能力。
特别是 \ac{VLM} 通过联合建模视觉与文本信息，显著拓展了模型对真实世界场景的理解边界，并广泛应用于对话助手、自主智能体及内容生成等关键领域。
然而，随着模型更深地融入现实世界系统，其潜藏的安全风险亦愈发凸显。在实际部署中，\ac{VLM}通常并非孤立运行，而是作为更大系统中的感知与决策模块。一旦模型在跨模态理解或安全判断上出现偏差，该错误可能通过后续模块被放大，最终导致系统级风险。
本文聚焦于这一风险链条的源头——模型自身的信息可靠性问题与输出安全。
当前 \ac{VLM} 面临的安全威胁已不再局限于单一模态：攻击者可通过文本提示诱导有害输出（文本风险），利用对抗性图像或不当视觉内容触发违规响应（图片风险），更可通过图文协同方式精心构造实施跨模态攻击（图文联合风险）。
此类风险的应对面临双重根本性挑战：
其一，安全目标与模型内生的指令遵循能力存在内在冲突——输出安全本质是要求模型拒绝风险内容，而通用能力则要求其忠实执行用户指令，这种张力导致模型常陷入"过度拒绝"与"拒绝不足"的帕累托困境；
其二，安全边界的主观性与不确定性使得第三方预定义规则难以普适，不同文化背景与应用场景对"有害"的界定存在显著差异，传统基于关键词的前置过滤方法往往泛化能力不足。
当前安全增强方法主要分为两条技术路线，但均存在显著瓶颈。
在推理与部署阶段，主流方案依赖输入输出审核\supercite{wang-adashield-eccv-2024,guan-deliberative-alignment-2024}（如Llama-Guard、Shield-Gemma等轻量分类器）
或推理时控制（如流式安全解码、Best-Of-N采样）。
尽管此类方法无需修改模型参数、部署灵活，但其防护效果高度依赖外部审查链路的完整性，且在面对开源模型的直接访问时完全失效；轻量审核模型的泛化能力亦难以应对频繁变化的安全规则与隐蔽的跨模态攻击模式。
在模型训练阶段，现有研究多局限于传统的指令微调（SFT）与直接偏好优化（DPO），通过"问题-拒答"对或人类反馈强化学习（RLHF）进行对齐。
然而，这种范式往往导致模型仅习得机械式拒绝的表面行为，而缺乏对"为何不安全"的深层认知，即存在显著的安全推理鸿沟（Safety Reasoning Gap）。
此外，现有方案多聚焦于单一维度的安全能力增强，未能有效整合外部知识验证与内部安全反思，难以应对知识密集型请求中由幻觉引发的安全风险。
针对上述局限，本文旨在构建一个系统化的多模态安全模型训练范式，重点解决以下核心问题：
（1）如何弥合安全推理鸿沟，使模型具备面对风险问题攻击时的自我反思与动态纠错能力？
（2）如何通过可信信息检索机制，应对知识密集型场景中的幻觉风险？
为此，本文提出融合推理时安全反思与训练时推理增强的创新框架。
具体而言，本文引入了多模态安全反思机制：模型在生成回复前，需显式分析图文上下文中的潜在恶意意图，对自身输出进行安全性评估与迭代修正，这不仅对应推理阶段的自反思控制，更通过安全思维链（Safety Chain-of-Thought）训练内化为模型的固有能力。
其次，我们设计了Agentic检索增强生成（RAG）框架，使模型能够自主决策是否调用外部知识库进行事实核查，从而在保证信息真实性的同时，从源头遏制基于错误知识的诱导攻击，尤其适用于时效性强、知识密集的安全场景。
本文的主要贡献包括：（1）提出并实现了从后训练阶段介入的多模态安全模型系统构建方案，突破传统"打补丁"式对齐的局限；（2）创新性地将Agentic RAG与安全反思机制深度融合，分别针对信息可靠性与生成安全性构建双重防线，显著提升了模型对复杂多模态攻击的鲁棒性；（3）通过系统实验验证了所提出范式在降低拒绝不足率、抵御越狱攻击及缓解幻觉风险方面的有效性。
本文后续结构如下：第二部分详细综述多模态安全风险与防御方法；第三部分阐述融合安全反思与Agentic RAG的训练数据构建及模型训练方法；第四部分介绍实验设置与结果分析；第五部分讨论当前局限与未来展望。

% vim:ts=4:sw=4